% Introduction paragraph for DDI Knowledge Graph project
% Covers: severity prediction, knowledge graph construction, DDI risk analysis, 
% alternative recommendations, and web application

\documentclass[11pt,a4paper]{article}
\usepackage[utf8]{inputenc}
\usepackage{geometry}
\geometry{margin=1in}

\title{\textbf{Introduction Paragraph for DDI Knowledge Graph Project}}
\author{}
\date{}

\begin{document}
\maketitle

\section*{Introduction}

While comprehensive drug-drug interaction (DDI) databases such as DrugBank catalog hundreds of thousands of interaction pairs, translating this information into actionable clinical guidance remains challenging: severity labels often lack evidence-based calibration, DDI data is siloed from mechanistic knowledge such as shared protein targets and metabolic pathways, and identifying therapeutically equivalent alternatives with lower interaction burden requires considerable manual effort. Given the complexity of analyzing the entire DrugBank database, this work focuses on cardiovascular and antithrombotic drugs---a therapeutic domain where polypharmacy is common and DDI consequences are particularly severe. We present an integrated pipeline comprising: (1) a severity recalibration methodology comparing zero-shot transformer-based natural language inference (BART-MNLI) against statistically-derived keyword classifiers, with keyword weights computed from log-likelihood ratios on DDInter-validated training data and classification thresholds optimized via percentile-based grid search; (2) a multi-relational biomedical knowledge graph integrating drugs, proteins, pathways, diseases, side effects, and therapeutic categories from DrugBank, SIDER, and CTD; (3) graph-based polypharmacy risk assessment with ATC classification-guided alternative drug recommendations; and (4) a web application integrating LLM capabilities across three core functions: text-based drug extraction from clinical notes, human-readable response generation for interaction reports, and a knowledge graph-aware chatbot that provides contextualized guidance based on the generated reports---enabling clinicians to rapidly assess cardiovascular medication regimens and identify safer therapeutic alternatives.

\end{document}
